\documentclass[10pt,a4paper]{amsart}
\usepackage[margin=19mm]{geometry}
\usepackage{amsmath,amssymb,mathtools}
\usepackage[scr=rsfs,cal=euler]{mathalfa}
\usepackage{xfrac}
\usepackage[unicode,hidelinks,bookmarks=false]{hyperref}
\newcommand{\ie}{i.e.\ }
\newcommand{\cf}{cf.\ }
\newcommand{\resp}{respectively\ }
\newcommand{\Subsection}{\subsection}
\providecommand{\nobreakdash}{\nobreak\hbox{-}}
\setlength{\emergencystretch}{2em}
\multlinegap0pt
\usepackage{xfrac}
\usepackage{nicefrac}
\usepackage{amssymb}
\usepackage{mathtools}
\usepackage[shortlabels]{enumitem}
\usepackage{booktabs}
\usepackage{xcolor}
\usepackage{float}
\usepackage{tikz}
\newtheorem{proposition}{Proposition}
\newcommand{\Cay}{\operatorname{Cay}}
\newcommand{\Z}{\mathbb{Z}}
\newcommand{\gl}{\operatorname{SDP_{GL}}}
\title{\textbf{Integrality Gap Bounds for the Goemans–Linial SDP on Finite Abelian Cayley Graphs}}
\author{Georgios Stamoulis}
\date{Source: arXiv:2609.05368v1 (CC BY 4.0)}
\begin{document}
\maketitle

\noindent\emph{Source and licence.} \textbf{Integrality Gap Bounds for the Goemans–Linial SDP on Finite Abelian Cayley Graphs}. Version arXiv:2609.05368v1 (CC BY 4.0), distributed by arXiv under the Creative Commons Attribution 4.0 International licence, http://creativecommons.org/licenses/by/4.0/, as recorded in the arXiv licence metadata for this exact version and shown on the abstract page https://arxiv.org/abs/2609.05368v1. Full licence notice: \texttt{licenses/arxiv-2609.05368-CC-BY-4.0.txt}.\par\smallskip

\noindent\textit{Editorial scope.} Counted unit: the article's proposition prop:cycle-exact (source0.tex lines 1207--1220). The article's complete original proof of it is delivered with it (source0.tex lines 1222--1245, one \texttt{proof} environment of 24 lines). No mathematics is written, repaired or re-derived: the statement and the proof are the article's own lines, in the article's own notation, and no retained line is substituted or re-flowed.  Retained uncounted as the source-local context the proof needs: lines 903--912.  Nothing was omitted from any retained span, so the package carries no omission record.  No retained line contains a citation, so the wrapper carries no bibliography and no bibliography entry.  No retained line contains a figure environment, an image-inclusion command or drawing code, so no image is delivered and the package declares no asset dependency.  Editorial formatting only, in addition: an AMS \texttt{amsart} preamble that restates the article's own declarations and macros used by the retained lines, namely the environments \texttt{proposition} on separate counters, together with the standard packages those lines need.  The retained environments print in this document as Proposition 1 in order of appearance, whereas the article prints the same items with the numbering of its own full document.  The complete native source of the article as distributed in the arXiv e-print for this exact version is bundled unchanged as \texttt{sources/209431/source0.tex}.
\begin{proposition}\label{prop:metric-averaging}
For each metric $d$ on $\Gamma$,  $N_{G^\sharp}(d)  ~\leq~ \rho(S) \cdot N_G(d)$
%\begin{equation}\label{eq:averaging-N}
%  N_{G^\sharp}(d)  ~\leq~ \rho(S) \cdot N_G(d).
%\end{equation}
and, consequently, $\gl(G^\sharp) \leq \rho(S) \cdot \gl(G)$.
%\begin{equation}\label{eq:averaging-arv}
%  \gl(G^\sharp) ~\leq~ \rho(S) \cdot \gl(G).
%\end{equation}
\end{proposition}

\begin{proposition}
% the Exact GL value on cycles
\label{prop:cycle-exact}
For each $n\geq 3$,
\begin{equation}\label{eq:cycle-exact}
  \gl(C_n)=\psi(C_n)
  =
  \begin{cases}
    4/n,& n \text{ is even},\\[6pt]
    (4n)/(n^2-1),&n\text{ is odd}.
  \end{cases}
\end{equation}
In fact, the same value is obtained if the SDP is relaxed further to all metrics satisfying the triangle inequalities, without the squared-Euclidean requirement.
\end{proposition}

\begin{proof}
Let $d$ be any metric on $\Z_n$. Averaging the generator $1$ makes the random step uniform over $\Z_n$ so $C_n^\sharp=\Cay(\Z_n,\Z_n)$ in normalized random-walk form where the zero step is just a self loop. Since $d(x,x)=0$ we have that
\[
  N_{C_n^\sharp}(d)
  =\frac1{n^2}\sum_{x \neq y}d(x,y) ~=~ D(d).
\]
By Proposition \ref{prop:metric-averaging} we have $D(d) \leq \alpha(n) \cdot N_{C_n}(d)$ and so every metric satisfies
\begin{equation}\label{eq:cycle-metric-lower}
  \frac{N_{C_n}(d)}{D(d)}
  ~\geq~\frac{1}{\alpha(n)}.
\end{equation}

Now let $m=\lfloor n/2\rfloor$ and take the interval cut $A=\{0,1,\ldots,m-1\}$. Exactly two undirected cycle edges cross the cut, equivalently four of the $2n$ directed transitions cross, so $N_{C_n}(A)=\nicefrac{2}{n}$. Its denominator is $D(A) = 2 \cdot \frac{m}{n} \left(1-\frac{m}{n}\right) = \frac{2m(n-m)}{n^2}$.
Therefore
\[
  \psi_{C_n}(A)=\frac{n}{m(n-m)}
  =
  \begin{cases}
    4/n,&n\text{ even},\\[3pt]
    4n/(n^2-1),&n\text{ odd}.
  \end{cases}
\]
Using the closed form for $\alpha$ the lower bound in \eqref{eq:cycle-metric-lower} is exactly the same quantity. Hence the metric relaxation, the GL relaxation, and the cut optimum all coincide.
\end{proof}

\end{document}
